\subsection{Hard patterns}

The aim of this section is to prove that for every $\ell \in [8104]$, {\sc Planar $\mathcal{C}_\ell$-Steiner Network} 
is $\mathsf{W}[1]$-hard parameterized by the number $k$ of terminals and does not admit a
$f(k) \cdot n^{o(k)}$ algorithm for any computable function $f$, unless $\mathsf{ETH}$ fails.
For each $\ell \in [8104]$, we give a reduction which transforms an instance of {\sc $k \times k$-Grid Tiling} (see \Cref{sec:diamonds} for a definition of this problem) into an instance of (edge-weighted)\footnote{We then use the same trick as in \Cref{sec:diamonds} to get rid of the edge-weights.} {\sc Planar $\mathcal{C}_\ell$-Steiner Network} with $O(k)$ terminals.
The constructed instances in each case are very similar and are based on a construction developed in the proof of \cite[Theorem 1.4]{ChitnisFHM20}, which we describe below as the \emph{main gadget}. We then show how to built upon this construction to handle the hard matching patterns (see \Cref{sec:orderedtough}) and the hard biclique patterns (see \Cref{sec:biclique}).\\

\noindent
\textbf{Main Gadget.} As mentioned above, we use the same construction as in the proof of \cite[Theorem 1.4]{ChitnisFHM20}. More precisely, given an integer $n > 0$ and a subset $S \subseteq [n] \times [n]$\footnote{Recall that, by assumption, $1 < x,y< n$ holds for every $(x,y) \in S$.}, we first construct an edge-weighted planar digraph $G(S)$ as follows. The graph $G(S)$ consists of an $n \times n$ grid where the horizontal edges are oriented towards the right and the vertical edges are oriented towards the bottom, that is, denoting by $x_{i,j}$ the vertex lying at the intersection of column $i$ and row $j$, there is an edge 
\begin{itemize}
\item $(x_{i+1,j},x_{i,j})$ (of weight 2) for every $i \in [n-1]$ and $j \in [n]$, and 
\item $(x_{i,j},x_{i,j+1})$ (of weight 2) for every $i \in [n]$ and $j \in [n-1]$. 
\end{itemize}
Then for every $(i,j) \in S$, we subdivide the edge $(x_{i-1,j},x_{i,j})$, by adding a vertex $y_{i,j}$ and the edges $(x_{i-1,j},y_{i,j})$ and $(y_{i,j},x_{i,j})$ (both of weight 1), and further add the edge $(x_{i,j-1},y_{i,j})$ (of weight 1). 
This concludes the construction of $G(S)$.
In the following, we call the vertices $x_{1,1},x_{1,2},\ldots,x_{1,n}$ the \emph{left vertices}, the vertices $x_{n,1}, x_{n,2}\ldots,x_{n,n}$ the \emph{right vertices}, the vertices $x_{1,n},x_{2,n}, \ldots,x_{n,n}$ the \emph{top vertices} and the vertices $x_{1,1},x_{2,1},\ldots, x_{n,1}$ the \emph{bottom vertices}. \\

Now given a collection $\mathcal{S} = \{S_{i,j}~|~i,j \in [k]\}$ of $k^2$ subsets of $[n] \times [n]$, the \emph{main gadget $MG(\mathcal{S})$} for $\mathcal{S}$ is constructed as follows.
\begin{itemize}
\item For every set $S_{i,j} \in \mathcal{S}$, we introduce a copy of the graph $G(S_{i,j})$ as constructed above. For every $i \in [n]$ and $j \in [n-1]$, the graph $G(S_{i,j+1})$ lies below the $G(S_{i,j})$; we add an edge (of weight 2) from each top vertex of $G(S_{i,j+1})$ to the bottom vertex of $G(S_{i,j})$ of the same index. Similarly, for every $i \in [n-1]$ and $j \in [n]$, the graph $G(S_{i,j})$ lies to the left of the graph $G(S_{i+1,j})$; we add an edge (of weight 2) from each right vertex of $G(S_{i,j})$ to the left vertex of $G(S_{i+1,j})$ of the same index.
\item We introduce $4k$ additional vertices $a_1,\ldots,a_k,b_1,\ldots,b_k,c_1,\ldots,c_k,d_1,\ldots,d_k$ and the following edges (we fix $\Delta = 5n^2$).
\begin{itemize}
\item For every $j \in [k]$ and every $i \in [n]$, there is an edge from $a_j$ to the left vertex of $G(S_{1,j})$ of index $i$ of weight $\Delta(n+1-i)$.
\item For every $j \in [k]$ and every $i \in [n]$, there is an edge from the right vertex of $G(S_{n,j})$ of index $i$ to $b_j$ of weight $\Delta i$.
\item For every $j \in [k]$ and every $i \in [n]$, there is an edge from $c_j$ to the top vertex of $G(S_{j,1})$ of index $i$ of weight $\Delta(n+1-i)$.
\item For every $j \in [k]$ and every $i \in [n]$, there is an edge from the bottom vertex of $G(S_{j,n})$ of index $i$ to $d_j$ of weight $\Delta i$.
\end{itemize}
\end{itemize}
This concludes the construction of the main gadget $MG(\mathcal{S})$ for $\mathcal{S}$ (see \Cref{fig:bimaingad} for an illustration of the main gadget $MG(\mathcal{S})$ with $n=4$ for $\mathcal{S} = \{S_{i,j}~|~i,j \in [3]\}$ where $S_{1,1} =\{(2,2),(2,3),(3,3)\}$ and $S_{i,j} = \emptyset$ for every $(i,j) \neq (1,1)$). We now set $I = \{a_i,b_i,c_i,d_i~|~i \in [k]\}$ and let $M$ be the induced matching $\{(a_i,b_i),(c_i,d_i)~|~i\in [k]\}$. We further set
\[
B^* = 2k(\Delta(n+1) + 2(k+1) +2k(n-1)).
\]

\begin{lemma}[\cite{ChitnisFHM20}]
\label{lem:grid-dsn}
The {\sc $k \times k$-Grid Tiling} instance $(k,n,\mathcal{S})$ has a solution if and only if the (edge-weighted) {\sc Planar $M$-Steiner Network} instance $(MG(\mathcal{S}),I,M)$ has a solution of weight at most $B^*-k^2$.
\end{lemma}

\begin{figure}
\centering
\begin{tikzpicture}[scale=.7]
\foreach \x in {0,5,10}
\foreach \y in {0,5,10}
{\pgfmathtruncatemacro{\a}{\x+1}
\pgfmathtruncatemacro{\e}{\x+3}
\pgfmathtruncatemacro{\b}{\x+4}
\pgfmathtruncatemacro{\c}{\y+1}
\pgfmathtruncatemacro{\f}{\y+2}
\pgfmathtruncatemacro{\d}{\y+4}

\pgfmathsetmacro{\u}{\x+.67}
\pgfmathsetmacro{\v}{\x+4.33}
\pgfmathsetmacro{\w}{\y+.67}
\pgfmathsetmacro{\z}{\y+4.33}
\filldraw[fill=lightgray] (\u,\w) rectangle (\v,\z);
\ifthenelse{\x=0}{\pgfmathtruncatemacro{\s}{\x+1}}{\ifthenelse{\x=5}{\pgfmathtruncatemacro{\s}{\x-3}}{\pgfmathtruncatemacro{\s}{\x-7}}}
\ifthenelse{\y=0}{\pgfmathtruncatemacro{\t}{\y+3}}{\ifthenelse{\y=5}{\pgfmathtruncatemacro{\t}{\y-3}}{\pgfmathtruncatemacro{\t}{\y-9}}}
\node[draw=none,label={[label distance=-.35cm]10:\footnotesize $G(S_{\s,\t})$}] at (\v,\z) {};


\foreach \i in {\c,...,\d}
\foreach \j in {\a,...,\b}
{
\node[mcirc] (\j\i) at (\j,\i) {};
}

\ifthenelse{\x=0}{
\pgfmathsetmacro{\r}{\y+2.5}
\node[circ,label=left:{\small $a_\t$}] (a\t) at (-.4,\r) {};
\foreach \i in {\c,...,\d}
{\pgfmathtruncatemacro{\o}{5-\i+\y}
\ifthenelse{\o=1}{
\draw[->,thick,>=stealth] (a\t) -- (\a\i) node[pos=.86,left] {{\fontsize{4}{10}\selectfont $\Delta$}};}
{\ifthenelse{\o=4}{\draw[->,thick,>=stealth] (a\t) -- (\a\i) node[pos=.33,right] {{\fontsize{4}{10}\selectfont $\o \Delta$}};}{\draw[->,thick,>=stealth] (a\t) -- (\a\i) node[pos=.63,above] {{\fontsize{4}{10}\selectfont $\o \Delta$}};}}}}{}

\ifthenelse{\x=10}{
\pgfmathsetmacro{\r}{\y+2.5}
\node[circ,label=right:{\small $b_\t$}] (b\t) at (15.4,\r) {};
\foreach \i in {\c,...,\d}
{\pgfmathtruncatemacro{\o}{\i-\y}
\ifthenelse{\o=1}{
\draw[->,thick,>=stealth] (\b\i) -- (b\t) node[pos=.59,left] {{\fontsize{4}{10}\selectfont $\Delta$}};}
{\ifthenelse{\o=4}{\draw[->,thick,>=stealth] (\b\i) -- (b\t) node[pos=.11,right] {{\fontsize{4}{10}\selectfont $\o \Delta$}};}{\draw[->,thick,>=stealth] (\b\i) -- (b\t) node[pos=.4,above] {{\fontsize{4}{10}\selectfont $\o \Delta$}};}}}}{}

\ifthenelse{\y=0}{
\pgfmathsetmacro{\r}{\x+2.5}
\node[circ,label=below:{\small $d_\s$}] (d\s) at (\r,-.4) {};
\foreach \j in {\a,...,\b}
{\pgfmathtruncatemacro{\o}{\j-\x}
\ifthenelse{\o=1}{
\draw[->,thick,>=stealth] (\j\c) -- (d\s) node[pos=.4,right] {{\fontsize{4}{10}\selectfont $\Delta$}};}
{\ifthenelse{\o=4}{\draw[->,thick,>=stealth] (\j\c) -- (d\s) node[pos=.4,label={[label distance=-.2cm]0:{\fontsize{4}{10}\selectfont $\o \Delta$}}] {};}{\draw[->,thick,>=stealth] (\j\c) -- (d\s) node[pos=.4,label={[label distance=-.2cm]0:{\fontsize{4}{10}\selectfont $\o \Delta$}}] {};}}}}{}

\ifthenelse{\y=10}{
\pgfmathsetmacro{\r}{\x+2.5}
\node[circ,label=above:{\small $c_\s$}] (c\s) at (\r,15.4) {};
\foreach \j in {\a,...,\b}
{\pgfmathtruncatemacro{\o}{5-\j+\x}
\ifthenelse{\o=1}{
\draw[->,thick,>=stealth] (c\s) -- (\j\d) node[pos=.6,right] {{\fontsize{4}{10}\selectfont $\Delta$}};}
{\ifthenelse{\o=4}{\draw[->,thick,>=stealth] (c\s) -- (\j\d) node[pos=.6,label={[label distance=-.2cm]0:{\fontsize{4}{10}\selectfont $\o \Delta$}}] {};}{\draw[->,thick,>=stealth] (c\s) -- (\j\d) node[pos=.6,label={[label distance=-.2cm]0:{\fontsize{4}{10}\selectfont $\o \Delta$}}] {};}}}}{}

\foreach \i in {\c,...,\d}
\foreach \j in {\a,...,\e}
{
\pgfmathtruncatemacro{\k}{\j+1}
\draw[->,thick,>=stealth] (\j\i) -- (\k\i);
}

\foreach \j in {\a,...,\b}
\foreach \i in {\d,...,\f}
{
\pgfmathtruncatemacro{\k}{\i-1}
\draw[->,thick,>=stealth] (\j\i) -- (\j\k);
}

\foreach \j in {\a,...,\b}
{\ifthenelse{\y=0}{}{
\pgfmathtruncatemacro{\k}{\y-1}
\draw[->,thick,>=stealth] (\j\c) -- (\j\k);}
}

\foreach \i in {\c,...,\d}
{\ifthenelse{\x=0}{}{
\pgfmathtruncatemacro{\k}{\x-1}
\draw[->,thick,>=stealth] (\k\i) -- (\a\i);}
}
}

\node[circr] (h) at (2,2.6) {};
\node[circr] (g) at (2,3.6) {};
\node[circr] (q) at (3,3.6) {};
\draw[->,thick,>=stealth,red] (12) -- (h);
\draw[->,thick,>=stealth,red] (13) -- (g);
\draw[->,thick,>=stealth,red] (23) -- (q);
\end{tikzpicture}
\caption{The main gadget $MG(\mathcal{S})$ with $n=4$ for $\mathcal{S} = \{S_{i,j}~|~i,j \in [3]\}$ where $S_{1,1} =\{(2,2),(2,3),(3,3)\}$ and $S_{i,j} = \emptyset$ for every $(i,j) \neq (1,1)$.}
\label{fig:bimaingad}
\end{figure}

\subsubsection{Hard matching patterns}
\label{sec:orderedtough}  

The aim of this section is to prove hardness for the hard matching patterns. We restate here the definition of these graphs for the reader's convenience.

\begingroup
\def\thetheorem{\ref{def:cleanedorderedtoughpair}}
\begin{definition}[$t$-hard matching pattern]
A \emph{$t$-hard matching pattern} is
an (acyclic) digraph $G$ constructed the following way. We start with disjoint vertex sets $A= \{a_1,\ldots,a_t\}$, $B = \{b_1,\ldots,b_t\}$, $C = \{c_1,\ldots,c_t\}$ and $D = \{d_1,\ldots,d_t\}$ and introduce the edges $(a_i,b_i)$ and $(c_i,d_i)$ for every $i\in[t]$.
Furthermore, we introduce into $G$ any combination of the following items:
\begin{enumerate}
\item either the directed path $a_1 \rightarrow a_2 \rightarrow \ldots \rightarrow a_t \rightarrow d_1 \rightarrow d_2 \rightarrow \ldots \rightarrow d_t$, or any of the directed paths $a_1 \rightarrow a_2 \rightarrow \ldots \rightarrow a_t$ and 
$d_1 \rightarrow d_2 \rightarrow \ldots \rightarrow d_t$;
\item either the directed path $c_1 \rightarrow c_2 \rightarrow \ldots \rightarrow c_t \rightarrow b_1 \rightarrow b_2 \rightarrow \ldots \rightarrow b_t$, or any of the directed paths $b_1 \rightarrow b_2 \rightarrow \ldots \rightarrow b_t$ and $c_1 \rightarrow c_2 \rightarrow \ldots \rightarrow c_t$;
\item an $S$-source for exactly one $S \in \{A,B,C,D,A \cup C,B \cup D,B \cup C,A \cup D\}$;
\item an $S$-sink for exactly one $S \in \{A,B,C,D,A \cup C,B \cup D,B \cup C,A \cup d\}$;
\item a vertex $r_{AD}$ such that $N^-(r_{AD}) = A$ and $N^+(r_{AD}) = D$;
\item a vertex $r_{CB}$ such that $N^-(r_{CB}) = Y$ and $N^+(r_{CB}) = B$.
\end{enumerate}
In particular, there are $5 \cdot 5 \cdot 9 \cdot 9 \cdot 2 \cdot 2$ types of $t$-hard matching patterns: we let $\mathcal{C}_5, \ldots, \mathcal{C}_{8104}$ be the 8100 classes that each contain all the $t$-hard matching patterns of a specific type for every~$t$.
\end{definition}
\addtocounter{theorem}{-1}
\endgroup

Formally, we aim to prove the following.

\begin{lemma}
\label{lem:cleanedtoughpair}
For every $\ell \in [5,8104]$, 
{\sc Planar $\mathcal{C}_\ell$-Steiner Network} is $\mathsf{W}[1]$-hard parameterized by the number $k$ of terminals and does not admit a $f(k) \cdot n^{o(k)}$ algorithm for any computable function $f$,
unless $\mathsf{ETH}$ fails.
\end{lemma}

To prove the lemma, we use the construction described above and further add vertices and edges so as to take into account the specific combination of the items in \Cref{def:cleanedorderedtoughpair} for each fixed $\ell \in [5,8104]$.\\

\noindent
\textbf{Reduction.} Consider $\ell \in [7,2031]$. Given an instance $(k,n,\mathcal{S}=\{S_{i,j}~|~i,j \in [k]\})$ of {\sc $k \times k$-Grid Tiling}, we construct an equivalent instance $(G(\mathcal{S}),T,D)$ of (edge-weighted) {\sc Planar $\mathcal{C}_\ell$-Steiner Network} as follows. Let $MG(\mathcal{S})$ be the main gadget for $\mathcal{S}$ as constructed above.
Then $T$ contains $\{a_i,b_i,c_i,d_i~|~i \in [k]\}$ (and possibly, as described below, some additional vertices depending on the class $\mathcal{C}\ell$) where $\{(a_i,b_i)~|~i \in [k]\}$ and $\{(c_i,d_i)~|~i \in [k]\}$ are the two perfect matchings contained in $D$. The graph $G(\mathcal{S})$ is then obtained from $MG(\mathcal{S})$ by introducing the directed paths
\[
a_1 \rightarrow \ldots \rightarrow a_k, d_1 \rightarrow \ldots \rightarrow d_k, c_1 \rightarrow \ldots \rightarrow c_k 
\text{ and } b_1 \rightarrow \ldots \rightarrow b_k
\]
where the weight of each newly added edge is set to 0. Furthermore, if the patterns in the class $C_\ell$ contain
\begin{itemize}
\item a source (that is, a vertex of item 3 in \Cref{def:cleanedorderedtoughpair}) then we add a vertex $s$ to $G(\mathcal{S})$ and $T$, and the edges $(s,a_1)$ and $(s,c_1)$, both of weight 0;
\item a sink (that is, a vertex of item 4 in \Cref{def:cleanedorderedtoughpair}) then we add a vertex $t$ to $G(\mathcal{S})$ and $T$, and the edges $(b_k,t)$ and $(d_k,t)$, both of weight 0;
\item the vertex of item 5 in \Cref{def:cleanedorderedtoughpair}, then we add a vertex $r_{AD}$ to $G(\mathcal{S})$ and $T$, and the edges $(a_k,r_{AD})$ and $(r_{AD},d_1)$, both of weight 0;
\item the vertex of item 6 in \Cref{def:cleanedorderedtoughpair}, then we add a vertex $r_{CB}$ to $G(\mathcal{S})$ and $T$, and the edges $(c_k,r_{CB})$ and $(r_{CB},b_1)$, both of weight 0.
\end{itemize}
This concludes the construction of $G(\mathcal{S})$. We let $D$ be the corresponding $k$-hard matching pattern of $\mathcal{C}_\ell$ on vertex set $T$. \Cref{lem:cleanedtoughpair} then follows from \Cref{lem:grid-dsn} and the lemma below.

\begin{lemma}
The (edge-weighted) {\sc Planar $\mathcal{C}_\ell$-Steiner Network} instance $(G(\mathcal{S}),T,D)$ has a solution of weight at most $B^*-k^2$ if and only if the (edge-weighted) {\sc Planar $M$-Steiner Network} instance $(MG(\mathcal{S}),I,M)$ has a solution of weight at most $B^*-k^2$.
\end{lemma}

\begin{proof}
If $E$ is a solution of $(MG(\mathcal{S}),I,M)$ of weight at most $B^*-k^2$ then it is easy to see that $E \cup (E(G(\mathcal{S})) \setminus E(MG(\mathcal{S})))$ is a solution of $(G(\mathcal{S}),T,D)$ of weight at most $B^*-k^2$.
Conversely, if $E$ is a solution of $(G(\mathcal{S}),T,D)$ of weight at most $B^*-k^2$ then the restriction of $E$ to $MG(\mathcal{S})$ is readily seen to be a solution of $(MG(\mathcal{S}),I,M)$ of weight at most $B^*-k^2$.
\end{proof}

\subsubsection{Hard biclique patterns}
\label{sec:biclique}

The aim of this section is to prove hardness for the hard biclique patterns. We restate here the definition of these graphs for the reader's convenience.

\begingroup
\def\thetheorem{\ref{def:cleanedbiclique}}
\begin{definition}[$t$-hard biclique pattern]
A \emph{$t$-hard biclique pattern} is an (acyclic) digraph $D$ constructed the following way. We start with two disjoint sets $A$ and $B$ with  $|A| = |B| = t$ and introduce every edge from $A$ to $B$. Furthermore, we introduce into $D$ any combination of the following items (see \Cref{fig:cleanedbiclique}):
\begin{enumerate}
\item an $A$-source;
\item a $B$-sink.
\end{enumerate}
In particular, there are $2 \cdot 2$ types of $t$-hard biclique patterns: we let $\mathcal{C}_1,\ldots,\mathcal{C}_4$ be the 4 classes that each contain all the $t$-hard biclique patterns of a specific type for every $t$.
\end{definition}
\addtocounter{theorem}{-1}
\endgroup

Formally, we aim to prove the following.

\begin{lemma}
\label{lem:cleanedbiclique}
For every $\ell \in [4]$,
{\sc Planar $\mathcal{C}_\ell$-Steiner Network} is $\mathsf{W}[1]$-hard parameterized by the number $k$ of terminals and does not admit a $f(k) \cdot n^{o(k)}$ algorithm for any computable function $f$,
unless $\mathsf{ETH}$ fails.
\end{lemma}

We only formally prove the statement for the class of all hard biclique patterns containing no further vertices 
as it will become clear from the proof that to handle the classes of all hard biclique patterns containing
\begin{itemize}
\item the sink vertex, it suffices to add a vertex $t$ and an edge $(t_i,t)$ for each $i \in [2k+1]$ in the construction below; or
\item the source vertex, it suffices to add a vertex $s$ and an edge $(s,s_i)$ for each $i \in [2k+1]$ in the construction below.
\end{itemize} 
In the following, we assume that the class of all hard biclique patterns with no further vertices is $\mathcal{C}_1$.\\

\newcommand{\biclique}{{\sc Planar $\mathcal{C}_1$-Steiner Network}}

\noindent
\textbf{Reduction.} Given an instance $(k,n,\mathcal{S}=\{S_{i,j}~|~i,j \in [k]\})$ of {\sc $k \times k$-Grid Tiling}, we construct an equivalent instance $(G(\mathcal{S}),T,D)$ of (edge-weighted) \biclique\ as follows. We start by constructing an auxiliary planar digraph $H$ consisting of $2(2k+1)$ distinguished vertices $s_1,\ldots,s_{2k+1},t_1,\ldots,t_{2k+1}$, and $2(2k+1)$ edge-disjoint directed paths $P_{i,j}$ with $i \in [2k+1]$ and $j \in \{i,i+1\}$ (where indices are taken modulo $2k+1$ henceforth), defined as follows.
\begin{itemize}
\item For every $i \in [2k+1]$, $P_{i,i}= s_iu_i^1\ldots u_i^{2k-1}t_i$ is a directed path from $s_i$ to $t_i$ of length $2k$.
\item For every $i \in [2k+1]$, $P_{i,i+1}$ is the directed path $s_i u^1_{i-1}u^2_{i-2}\ldots u^j_{i-j}\ldots u^{2k-1}_{i+2}t_{i+1}$ from $s_i$ to $t_{i+1}$ of length $2k$.
\end{itemize}
(For $k=3$, the graph $H$ can be obtained from the graph depicted in \Cref{fig:bicliquered} by ignoring the blue vertices and contracting every grey box into a single vertex.) Note that by construction, the following holds.

\begin{observation}
\label{obs:uniqupath}
For every $i \in [2k+1]$, there is a unique path from $s_i$ to $t_i$ ($t_{i+1}$, respectively) in $H$, namely $P_{i,i}$ ($P_{i,i+1}$, respectively). 
\end{observation}

\noindent
The graph $G(\mathcal{S})$ is then obtained from $H$ as follows (see \Cref{fig:bicliquered}).
\begin{itemize}
\item We subdivide the edge $(s_{k+1},u^{k+1}_1)$ by adding the vertex $a_1$ and the edges $(s_{k+1},a_1)$ and $(a_1,u^{k+1}_1)$; and for every $2 \leq i \leq k$, we subdivide the edge $(u^i_{k+1-i},u^i_{k+2-i})$ by adding the vertex $a_{k+2-i}$ and the edges $(u^i_{k+1-i},a_{k+2-i})$ and $(a_{k+2-i},u^i_{k-i})$.
\item We subdivide the edge $(u^2_{2k-1},t_2)$ by adding the vertex $b_k$ and the edges $(u^2_{2k-1},b_k)$ and $(b_k,t_2)$; and for every $3 \leq i \leq k+1$, we subdivide the edge $(u^i_{2k+1-i},u^i_{2k+2-i})$ by adding the vertex $b_{k+2-i}$ and the edges $(u^i_{2k+1-i},b_{k+2-i})$ and $(b_{k+2-i},u^i_{2k+2-i})$.
\item We subdivide the edge $(s_{k+2},u^{k+1}_1)$ by adding the vertex $c_1$ and the edges $(s_{k+2},c_1)$ and $(c_1,u^{k+1}_1)$; and for every $k+3 \leq i \leq 2k+1$, we subdivide the edge $(u^{k+2}_{i-k-2},u^{k+1}_{i-k-2})$ by adding the vertex $c_{i-k-1}$ and the edges $(u^{k+2}_{i-k-2},c_{i-k-1})$ and $(c_{i-k-1},u^{k+1}_{i-k-2})$.
\item For every $k+2 \leq i \leq 2k+1$, we subdivide the edge $(u^2_{i-k-2},u^1_{i-k-2})$ by adding the vertex $d_{i-k-1}$ and the edges $(u^2_{i-k-2},d_{i-k-1})$ and $(d_{i-k-1},u^1_{i-k-2})$.
\item For every $1\leq i,j \leq k$, we replace the vertex $u^{j+1}_{k-j+i}$ with a copy of the graph $G(S_{i,k+1-j})$ and add the necessary edges so that the subgraph of $G$ induced by $\{a_i,b_i,c_i,d_i~|~ 1 \leq i \leq k\} \cup \bigcup_{1 \leq i,j \leq k} V(G(S_{i,j}))$ is isomorphic to $MG(\mathcal{S})$ (and has the same edge-weight function).
\end{itemize}
The weight of each edge outside the copy of the main gadget $MG(\mathcal{S})$ is set to 0. This concludes the construction of $G(\mathcal{S})$. We set $T = \{s_i,t_i~|~i\in [2k+1]\}$ and let the demand graph $D$ be the corresponding $(\{s_i~|~i \in [2k+1]\},\{t_i~|~i \in [2k+1]\})$-biclique. \Cref{lem:cleanedbiclique} for $\ell=1$ then follows from \Cref{lem:grid-dsn} and the lemma below.

\begin{figure}
\centering
\begin{tikzpicture}
\foreach \i in {1,...,7}
{\pgfmathsetmacro{\x}{\i-.75}
\node[circR,label={above:\small $s_\i$}] (a\i) at (\x,\x) {};
}

\foreach \i in {1,...,7}
{\pgfmathsetmacro{\y}{\i-.75}
\pgfmathsetmacro{\x}{\y+6}
\node[circR,label={below:\small $t_\i$}] (b\i) at (\x,\y) {};
}

\foreach \i in {2,...,6}
{\pgfmathsetmacro{\x}{\i+5.25}
\node[scirc] (\i7) at (\x,6.25) {};
}

\foreach \i in {1,...,6}
{\pgfmathtruncatemacro{\j}{\i+1}
\pgfmathsetmacro{\a}{\i-1}
\pgfmathsetmacro{\1}{\a*.5}
\pgfmathsetmacro{\b}{-\1}
\pgfmathsetmacro{\c}{\b+.25}
\pgfmathsetmacro{\2}{\1+7}
\pgfmathsetmacro{\3}{\a+.25}
\pgfmathsetmacro{\4}{\i+6.25}
\ifthenelse{\i=1 \OR \i=6}{\ifthenelse{\i=1}{\draw[->,thick,>=stealth] (a\i)-- (\b,\c) -- (\b,\2) -- (\4,\2) -- (\j7);}{\draw[->,thick,>=stealth] (a\i) -- (\3,\c) -- (\b,\c) -- (\b,\2) -- (\4,\2) -- (b\j);}}{\draw[->,thick,>=stealth] (a\i) -- (\3,\c) -- (\b,\c) -- (\b,\2) -- (\4,\2) -- (\j7);}
}

\foreach \i in {2,...,6}
{\ifthenelse{\i=6}{\draw[->,thick,>=stealth] (67) -- (b6);}{
\pgfmathtruncatemacro{\k}{\i+1}
\foreach \j in {6,...,\k}
{\pgfmathsetmacro{\x}{\i+5.25}
\pgfmathsetmacro{\2}{\j-.75}
\pgfmathtruncatemacro{\l}{\j+1}
\node[scirc] (\i\j) at (\x,\2) {};
\ifthenelse{\j=\k}{\draw[->,thick,>=stealth] (\i\l) -- (\i\j); \draw[->,thick,>=stealth] (\i\j) -- (b\i);}{\draw[->,thick,>=stealth] (\i\l) -- (\i\j);}
}
}
}

\node[scirc] (16) at (6.25,5.25) {};
\draw[->,thick,>=stealth] (a7) -- (16);
\node[scirc] (15) at (6.25,4.25) {};
\draw[->,thick,>=stealth] (16) -- (15);
\node[scirc] (05) at (5.25,4.25) {};
\draw[->,thick,>=stealth] (a6) -- (05);

\node[scirc] (-23) at (3.25,2.25) {};
\draw[->,thick,>=stealth] (a4) -- (-23);
\node[scirc] (-22) at (3.25,1.25) {};
\draw[->,thick,>=stealth] (-23) -- (-22);
\node[scirc] (-32) at (2.25,1.25) {};
\draw[->,thick,>=stealth] (a3) -- (-32);

\foreach \i in {1,...,5}
{\pgfmathsetmacro{\x}{\i+.25}
\pgfmathtruncatemacro{\j}{5-\i}
\node[scirc] (-\j1) at (\x,.25) {};
}

\draw[->,thick,>=stealth] (-22) -- (-21);
\draw[->,thick,>=stealth] (-32) -- (-31);
\draw[->,thick,>=stealth] (a2) -- (-41);

\foreach \i in {4,...,1}
{\pgfmathtruncatemacro{\j}{\i-1}
\draw[->,thick,>=stealth] (-\i1) -- (-\j1);
}
\draw[->,thick,>=stealth] (a1) -- (-41);
\draw[->,thick,>=stealth] (-01) -- (b1);

\draw[->,thick,>=stealth] (a2) -- (-32);
\draw[->,thick,>=stealth] (-32) -- (-22);

\draw[->,thick,>=stealth] (a3) -- (-23);

\draw[->,thick,>=stealth] (24) -- (34);

\foreach \i in {5,6,7}
{\pgfmathtruncatemacro{\k}{\i-5}
\pgfmathtruncatemacro{\l}{\i-2}
\foreach \j in {\k,...,\l}
{\pgfmathtruncatemacro{\p}{\j+1}
\ifthenelse{\j=\k}{\draw[->,thick,>=stealth] (a\i) -- (\j\i);\draw[->,thick,>=stealth] (\j\i) -- (\p\i);}{
\draw[->,thick,>=stealth] (\j\i) -- (\p\i);}
}
}

\draw[->,thick,>=stealth] (a7) -- (b1);
\draw[->,thick,>=stealth] (a2) -- (b2);


\foreach \i in {3,...,7}
{\pgfmathtruncatemacro{\j}{\i-1}
\draw[->,thick,>=stealth] (a\i) -- (\j\i);
\draw[->,thick,>=stealth] (\j\i) -- (b\i);
}

\draw[->,thick,>=stealth] (a4) -- (24);
\draw[->,thick,>=stealth] (a5) -- (-11);
\draw[->,thick,>=stealth] (a6) -- (-01);


\draw[->,thick,>=stealth] (a5) -- (4.25,3.5);
\draw[->,thick,>=stealth] (05) -- (5.25,3.5);
\draw[->,thick,>=stealth] (15) -- (6.25,3.5);

\draw[->,thick,>=stealth] (a4) -- (4,3.25); 
\draw[->,thick,>=stealth] (-23) -- (4,2.25);
\draw[->,thick,>=stealth] (-22) -- (4,1.25);

\foreach \i in {4,5,6}
\foreach \j in {1,2,3}
{\pgfmathsetmacro{\x}{\i+.5}
\pgfmathsetmacro{\y}{\j+.5}
\filldraw[fill=lightgray,very thick] (\i,\j) rectangle (\x,\y);
}

\foreach \i in {4,5}
\foreach \j in {1,2,3}
{\pgfmathsetmacro{\x}{\i+.5}
\pgfmathsetmacro{\u}{\i+1}
\pgfmathsetmacro{\y}{\j+.25}
\draw[->,thick,>=stealth] (\x,\y) -- (\u,\y);
}

\foreach \j in {3,2}
\foreach \i in {4,5,6}
{\pgfmathsetmacro{\x}{\i+.25}
\pgfmathsetmacro{\u}{\j-.5}
\draw[->,thick,>=stealth] (\x,\j) -- (\x,\u);
}

\foreach \i in {2,3,4}
{\pgfmathsetmacro{\y}{\i-.75}
\pgfmathtruncatemacro{\j}{5-\i}
\node[circb,label={[blue] below left:\tiny $a_\j$}] at (3.75,\y) {};
}

\foreach \i in {2,3,4}
{\pgfmathsetmacro{\y}{\i-.75}
\pgfmathtruncatemacro{\j}{5-\i}
\node[circb,label={[blue] above right:\tiny $b_\j$}] at (6.75,\y) {};
}

\foreach \i in {4,5,6}
{\pgfmathsetmacro{\x}{\i+.25}
\pgfmathtruncatemacro{\j}{\i-3}
\node[circb,label={[blue] below left:\tiny $d_\j$}] at (\x,.75) {};
}

\foreach \i in {4,5,6}
{\pgfmathsetmacro{\x}{\i+.25}
\pgfmathtruncatemacro{\j}{\i-3}
\node[circb,label={[blue] above right:\tiny $c_\j$}] at (\x,3.75) {};
}

\draw[dashed,blue,thick] (3.75,.75) rectangle (6.75,3.75);
\end{tikzpicture}
\caption{An illustration of the reduction from {\sc $k \times k$-Grid Tiling} to \biclique\ with $k=3$ (the red vertices are the terminals and the dashed blue square together with the grey boxes represent the main gadget).}
\label{fig:bicliquered}
\end{figure}

\begin{lemma}
The (edge-weighted) \biclique\ instance $(G(\mathcal{S}),T,D)$ has a solution of weight at most $B^*-k^2$ if and only if the (edge-weighted) {\sc Planar $M$-Steiner Network} instance $(MG(\mathcal{S}),I,M)$ has a solution of weight at most $B^*-k^2$.
\end{lemma}

\begin{proof}
If $E$ is a solution of $(MG(\mathcal{S}),I,M)$ of weight at most $B^*-k^2$ then it is not difficult to see that $E \cup (E(G(\mathcal{S})) \setminus E(MG(\mathcal{S})))$ is a solution of $(G(\mathcal{S}),T,D)$ of weight at most $B^*-k^2$. Conversely, let $E$ be a solution of $(G(\mathcal{S}),T,D)$ of weight at most $B^*-k^2$. Since for every $2 \leq i \leq k+1$, there is a unique path in $H$ from $s_i$ to $t_i$, namely $P_{i,i}$, and for every $k+2 \leq i \leq 2k+1$, there is a unique path in $H$ from $s_i$ to $t_{i+1}$, namely $P_{i,i+1}$, the restriction $E'$ of $E$ to $MG(\mathcal{S})$ is a solution of $(MG(\mathcal{S}),I,M)$; and since every edge in $E(G(\mathcal{S})) \setminus E(MG(\mathcal{S}))$ has weight 0, the weight of $E'$ is that of $E$, that is, $E'$ has weight at most $B^*-k^2$.
\end{proof}